% Distilled blueprint for Evans, Chapter 7 (Linear Evolution Equations).
% Formal declarations, metadata, and citation identifiers are preserved; exposition and exercises are omitted.

\section{Second-Order Parabolic Equations}
\label{sec:second-order-parabolic-equations}

\subsection{Definitions}

\subsubsection{Parabolic equations}

\begin{definition}[Uniform parabolicity]
\label{def:uniform-parabolicity}
\dcref{ch7:7.1.1-def-1}
\group{Parabolic PDE}
We say that the partial differential operator $\frac{\partial}{\partial t} + L$ is (uniformly) parabolic if there exists a constant $\theta > 0$ such that
\[
\sum_{i,j=1}^{n} a^{ij}(x,t)\xi_i \xi_j \geq \theta|\xi|^2
\tag{4}
\]
for all $(x,t) \in U_T$, $\xi \in \mathbb{R}^n$.
\end{definition}

\begin{remark}[Parabolic operators are elliptic at each fixed time]
\label{rem:parabolic-operator-elliptic-fixed-time}
\uses{def:uniform-parabolicity, def:uniform-ellipticity}
\dcref{ch7:7.1.1-rem-1}
\group{Parabolic PDE}

For every fixed $t\in[0,T]$, $L(t)$ is uniformly elliptic in the spatial variable with ellipticity constant $\theta$.
\end{remark}

\subsubsection{Weak solutions}

\begin{remark}[Weak parabolic time derivative]
\label{rem:motivation-weak-solution-parabolic}
\uses{def:sobolev-space-time-banach, def:dual-space-h-minus-1}
\dcref{ch7:7.1.1-rem-2}
\group{Parabolic PDE}

Regard $u$ and $f$ as maps $[0,T]\to H_0^1(U)$ and $[0,T]\to L^2(U)$. Formal spatial integration by parts gives
\[
(u',v)+B[u,v;t]=(f,v) \qquad (v\in H_0^1(U)).
\tag{9}
\]
Moreover,
\[
u_t=g^0+\sum_{j=1}^n g^j_{x_j},\qquad
g^0=f-\sum_{i=1}^n b^iu_{x_i}-cu,\quad
g^j=\sum_{i=1}^n a^{ij}u_{x_i},
\tag{10}
\]
and hence
\[
\|u_t\|_{H^{-1}(U)}\le C\bigl(\|u\|_{H_0^1(U)}+\|f\|_{L^2(U)}\bigr).
\]
Thus the weak identity is interpreted with $u'\in H^{-1}(U)$ and the dual pairing $\langle u',v\rangle$.
\end{remark}

\begin{definition}[Weak solution of the parabolic initial/boundary-value problem]
\label{def:weak-solution-parabolic-bvp}
\uses{rem:motivation-weak-solution-parabolic}
\dcref{ch7:7.1.1-def-2}
\group{Parabolic PDE}

We say a function
\[
\mathbf{u} \in L^2(0,T; H_0^1(U)), \text{ with } \mathbf{u}' \in L^2(0,T; H^{-1}(U)),
\]
is a weak solution of the parabolic initial/boundary-value problem (1) provided

\quad (i) $\langle \mathbf{u}', v \rangle + B[\mathbf{u},v;t] = (\mathbf{f},v)$

for each $v \in H_0^1(U)$ and a.e. time $0 \le t \le T$, and

\quad (ii) $\mathbf{u}(0) = g$.
\end{definition}

\begin{remark}[Continuity of the weak solution]
\label{rem:weak-solution-parabolic-continuity}
\uses{def:weak-solution-parabolic-bvp, thm:more-calculus-h1-hminus1}
\dcref{ch7:7.1.1-rem-3}
\group{Parabolic PDE}

Every weak solution belongs to $C([0,T];L^2(U))$, so its initial value is well-defined.
\end{remark}

\subsection{Existence of Weak Solutions}

\subsubsection{Galerkin approximations}

\begin{theorem}[Construction of approximate solutions]
\label{thm:parabolic-galerkin-approximate-solutions}
\dcref{ch7:7.1-thm-1}
\group{Parabolic PDE}
For each integer $m = 1, 2, \ldots$ there exists a unique function $u_m$ of the form (14) satisfying (15), (16).
\end{theorem}

\begin{proof}
\sketch Expand $u_m(t)$ in the first $m$ basis functions. The Galerkin identities form a linear finite-dimensional ODE system with the prescribed initial coefficients, so standard ODE theory gives a unique absolutely continuous solution.
\end{proof}

\subsubsection{Energy estimates}

\begin{theorem}[Energy estimates]
\label{thm:parabolic-energy-estimates}
\uses{thm:parabolic-galerkin-approximate-solutions, thm:energy-estimates}
\dcref{ch7:7.1-thm-2}
\group{Energy Estimates}

There exists a constant $C$, depending only on $U, T$ and the coefficients of $L$, such that
\[
\max_{0 \leq t \leq T} \|u_m(t)\|_{L^2(U)} + \|u_m\|_{L^2(0,T; H^1(U))} + \|u_m'\|_{L^2(0,T; H^{-1}(U))}
\]
\[
\leq C(\|f\|_{L^2(0,T; L^2(U))} + \|g\|_{L^2(U)})
\]
for $m = 1, 2, \ldots$ .
\end{theorem}

\begin{proof}
\sketch Test the Galerkin system with $u_m$, use ellipticity and Cauchy inequality to control lower-order terms, and apply Gronwall. Testing with $u_m^{\prime}$ and using the first estimate controls the time derivative and the terminal energy.
\end{proof}

\begin{theorem}[Existence of weak solution]
\label{thm:parabolic-existence-weak-solution}
\uses{thm:parabolic-galerkin-approximate-solutions, thm:parabolic-energy-estimates, def:weak-solution-parabolic-bvp}
\dcref{ch7:7.1-thm-3}
\group{Parabolic PDE}

There exists a weak solution of (11).
\end{theorem}

\begin{proof}
\sketch The uniform energy estimates give weakly convergent subsequences for $u_m$ and $u_m^{\prime}$. Pass to the limit in the Galerkin identities using density of the basis; integration by parts in time identifies the initial value and the abstract continuity theorem supplies the continuous $L^2$ representative.
\end{proof}

\begin{theorem}[Uniqueness of weak solutions]
\label{thm:parabolic-uniqueness-weak-solution}
\uses{def:weak-solution-parabolic-bvp, thm:energy-estimates}
\dcref{ch7:7.1-thm-4}
\group{Parabolic PDE}

A weak solution of (11) is unique.
\end{theorem}

\begin{proof}
\sketch Apply the energy identity to the difference of two solutions. Ellipticity, Cauchy inequality, and Gronwall force its $L^2$ norm to vanish.
\end{proof}

\subsection{Regularity}

\begin{remark}[Formal parabolic regularity estimates]
\label{rem:parabolic-regularity-motivation}
\dcref{ch7:7.1.3-rem-1}
\group{Regularity}
For a smooth rapidly decaying solution of $u_t-\Delta u=f$ on $\mathbb R^n\times(0,T)$ with $u(0)=g$, expanding $\|u_t-\Delta u\|_{L^2}^2$ and integrating by parts gives
\[
\sup_{0\le t\le T}\int_{\mathbb R^n}|Du|^2\,dx
+\int_0^T\int_{\mathbb R^n}u_t^2+|D^2u|^2\,dx\,dt
\le C\left(\int_0^T\int_{\mathbb R^n}f^2\,dx\,dt+\int_{\mathbb R^n}|Dg|^2\,dx\right).
\tag{39}
\]
Differentiating in time, applying the same energy estimate, and using $-\Delta u=f-u_t$ yields
\[
\sup_{0\le t\le T}\int_{\mathbb R^n}|u_t|^2+|D^2u|^2\,dx
+\int_0^T\int_{\mathbb R^n}|Du_t|^2\,dx\,dt
\le C\left(\int_0^T\int_{\mathbb R^n}f_t^2+f^2\,dx\,dt+\int_{\mathbb R^n}|D^2g|^2\,dx\right).
\tag{44}
\]
\end{remark}

\begin{theorem}[Improved regularity]
\label{thm:parabolic-improved-regularity}
\uses{def:weak-solution-parabolic-bvp, thm:boundary-h2-regularity}
\dcref{ch7:7.1-thm-5}
\group{Regularity}

(i) Assume
\[
g \in H_0^1(U), \quad f \in L^2(0,T; L^2(U)).
\]
Suppose also $u \in L^2(0,T; H_0^1(U))$, with $u' \in L^2(0,T; H^{-1}(U))$, is the weak solution of
\[
\begin{cases}
u_t + Lu = f & \text{in } U_T \\
u = 0 & \text{on } \partial U \times [0,T] \\
u = g & \text{on } U \times \{t = 0\}.
\end{cases}
\]
Then in fact
\[
u \in L^2(0,T; H^2(U)) \cap L^\infty(0,T; H_0^1(U)), \quad u' \in L^2(0,T; L^2(U)),
\]
and we have the estimate
\[
\operatorname{ess\,sup}_{0 \leq t \leq T} \|u(t)\|_{H_0^1(U)} + \|u\|_{L^2(0,T; H^2(U))} + \|u'\|_{L^2(0,T; L^2(U))}
\]
\[
\leq C \left( \|f\|_{L^2(0,T; L^2(U))} + \|g\|_{H_0^1(U)} \right),
\tag{46}
\]
the constant $C$ depending only on $U$, $T$ and the coefficients of $L$.

(ii) If, in addition,
\[
g \in H^2(U), \quad f' \in L^2(0,T;L^2(U)),
\]
then
\[
u \in L^\infty(0,T;H^2(U)), \quad u' \in L^\infty(0,T;L^2(U)) \cap L^2(0,T;H_0^1(U)),
\]
\[
u'' \in L^2(0,T;H^{-1}(U)),
\]
with the estimate
\[
\operatorname{ess\,sup}_{0 \leq t \leq T} \left( \|u(t)\|_{H^2(U)} + \|u'(t)\|_{L^2(U)} + \|u\|_{L^2(0,T;H_0^1(U))} \right.
\]
\[
\left. + \|u''\|_{L^2(0,T;H^{-1}(U))} \right) \leq C \left( \|f\|_{H^1(0,T;L^2(U))} + \|g\|_{H^2(U)} \right).
\tag{47}
\]
\end{theorem}

\begin{proof}
\sketch Differentiate the Galerkin equations in time and use the compatibility of the initial data to obtain uniform estimates for $u_m^{\prime}$ and $u_m^{\prime\prime}$. Pass to the limit, then apply elliptic boundary regularity to $Lu=f-u^{\prime}$ at almost every time.
\end{proof}

\begin{remark}[Precise versions of the formal estimates]
\label{rem:parabolic-improved-regularity-remark}
\uses{thm:parabolic-improved-regularity, rem:parabolic-regularity-motivation}
\dcref{ch7:7.1.3-rem-2}
\group{Regularity}

Parts (i) and (ii) of \cref{thm:parabolic-improved-regularity} are the variable-coefficient, bounded-domain versions of (39) and (44).
\end{remark}

\begin{remark}[Symmetric operators and eigenfunction bases]
\label{rem:parabolic-regularity-symmetric-eigenfunctions}
\uses{thm:parabolic-improved-regularity}
\dcref{ch7:7.1.3-rem-3}
\group{Regularity}

When $L$ is symmetric, the Galerkin basis may be chosen as an $H_0^1(U)$ eigenfunction basis of $L$, eliminating estimate (54).
\end{remark}

\begin{theorem}[Higher regularity]
\label{thm:parabolic-higher-regularity}
\uses{thm:parabolic-improved-regularity, thm:higher-boundary-regularity, thm:mappings-into-better-spaces}
\dcref{ch7:7.1-thm-6}
\group{Regularity}

Assume
\[
g \in H^{2m+1}(U), \quad \frac{d^k f}{dt^k} \in L^2(0,T; H^{2m-2k}(U)) \quad (k=0,\ldots,m).
\]
Suppose also the following $m^{\text{th}}$-order compatibility conditions hold:
\[
\begin{cases}
g_0 := g \in H^1_0(U), \quad g_1 := f(0) - Lg_0 \in H^1_0(U), \\
\ldots, \quad g_m := \dfrac{d^{m-1}}{dt^{m-1}}f(0) - Lg_{m-1} \in H^1_0(U).
\end{cases}
\]
Then
\[
\frac{d^k u}{dt^k} \in L^2(0,T; H^{2m+2-2k}(U)) \quad (k=0,\ldots,m+1);
\]
and we have the estimate
\[
\sum_{k=0}^{m+1} \left\| \frac{d^k u}{dt^k} \right\|_{L^2(0,T;H^{2m+2-2k}(U))}
\]
\[
\leq C \left( \sum_{k=0}^{m} \left\| \frac{d^k f}{dt^k} \right\|_{L^2(0,T;H^{2m-2k}(U))} + \|g\|_{H^{2m+1}(U)} \right),
\tag{55}
\]
the constant $C$ depending only on $m, U, T$ and the coefficients of $L$.
\end{theorem}

\begin{proof}
\sketch Induct on time derivatives. Differentiate the equation, apply the preceding parabolic estimate to each differentiated problem using the compatibility data, and recover spatial derivatives from higher elliptic regularity; the abstract time-continuity estimate gives the stated continuous-in-time spaces.
\end{proof}

\begin{remark}[Compatibility conditions]
\label{rem:parabolic-compatibility-conditions}
\uses{thm:parabolic-higher-regularity, thm:mappings-into-better-spaces}
\dcref{ch7:7.1.3-rem-4}
\group{Regularity}

The time traces satisfy
\[
f(0)\in H^{2m-1}(U),\quad f'(0)\in H^{2m-3}(U),\quad\ldots,\quad f^{(m-1)}(0)\in H^1(U),
\]
and consequently
\[
g_0\in H^{2m+1}(U),\quad g_1\in H^{2m-1}(U),\quad\ldots,\quad g_m\in H^1(U).
\]
Compatibility additionally requires every $g_j$ to have zero trace on $\partial U$.
\end{remark}

\begin{theorem}[Infinite differentiability]
\label{thm:parabolic-infinite-differentiability}
\uses{thm:parabolic-higher-regularity}
\dcref{ch7:7.1-thm-7}
\group{Regularity}

Assume
\[
g \in C^\infty(\bar{U}), \quad f \in C^\infty(\bar{U}_T),
\]
and the $m$-th order compatibility conditions hold for $m = 0, 1, \ldots$ . Then the parabolic initial/boundary-value problem (11) has a unique solution
\[
u \in C^\infty(\bar{U}_T).
\]
\end{theorem}

\begin{proof}
\sketch Iterate higher regularity for all orders and combine the resulting space-time Sobolev bounds with Sobolev embedding.
\end{proof}

\begin{remark}[Interior estimates]
\label{rem:parabolic-interior-estimates}
\uses{thm:parabolic-infinite-differentiability, thm:interior-h2-regularity}
\dcref{ch7:7.1.3-rem-5}
\group{Regularity}

Interior analogues of the global regularity estimates require no boundary compatibility conditions.
\end{remark}

\subsection{Maximum Principles}

\subsubsection{Weak maximum principle}

\begin{theorem}[Weak maximum principle]
\label{thm:parabolic-weak-maximum-principle}
\uses{def:uniform-parabolicity, thm:weak-maximum-principle-c-zero}
\dcref{ch7:7.1-thm-8}
\group{Maximum Principles}

Assume $u \in C^2_1(U_T) \cap C(\overline{U_T})$ and
\[
c = 0 \quad \text{in } U_T.
\tag{62}
\]
(i) If
\[
u_t + Lu \leq 0 \quad \text{in } U_T,
\tag{63}
\]
then
\[
\max_{\bar U_T} u = \max_{\Gamma_T} u.
\]
(ii) Likewise, if
\[
u_t + Lu \geq 0 \quad \text{in } U_T,
\tag{64}
\]
then
\[
\min_{\bar U_T} u = \min_{\Gamma_T} u.
\]
\end{theorem}

\begin{proof}
\sketch At a positive interior space-time maximum with positive time, the spatial Hessian is negative semidefinite and $u_t\ge0$, contradicting the strict differential inequality. Add a small exponential barrier to reduce the nonstrict case to the strict one and let the perturbation vanish.
\end{proof}

\begin{remark}[Subsolution and supersolution terminology]
\label{rem:parabolic-subsolution-supersolution}
\uses{thm:parabolic-weak-maximum-principle}
\dcref{ch7:7.1.4-rem-1}
\group{Maximum Principles}

Inequality (63) defines a subsolution, whose maximum lies on the parabolic boundary $\Gamma_T$; (64) defines a supersolution, whose minimum lies on $\Gamma_T$.
\end{remark}

\begin{theorem}[Weak maximum principle for $c \geq 0$]
\label{thm:parabolic-weak-maximum-principle-c-nonneg}
\uses{thm:parabolic-weak-maximum-principle}
\dcref{ch7:7.1-thm-9}
\group{Maximum Principles}

Assume $u \in C_1^2(U_T) \cap C(\bar{U}_T)$ and
\[
c \geq 0 \text{ in } U_T.
\]
(i) If
\[
u_t + Lu \leq 0 \text{ in } U_T,
\]
then
\[
\max_{\bar{U}_T} u \leq \max_{\Gamma_T} u^+.
\]
(ii) If
\[
u_t + Lu \geq 0 \text{ in } U_T,
\]
then
\[
\min_{\bar{U}_T} u \geq - \max_{\Gamma_T} u^-.
\]
\end{theorem}

\begin{proof}
\sketch On the positive set, remove the nonnegative zeroth-order term and apply the zero-order-free parabolic maximum principle; apply the result to $-u$ for the lower bound.
\end{proof}

\begin{remark}[Absolute value bound]
\label{rem:parabolic-weak-max-principle-abs-value}
\uses{thm:parabolic-weak-maximum-principle-c-nonneg}
\dcref{ch7:7.1.4-rem-2}
\group{Maximum Principles}

If $u_t+Lu=0$ in $U_T$, then
\[
\max_{\bar U_T}|u|=\max_{\Gamma_T}|u|.
\]
\end{remark}

\begin{remark}[Maximum principles with $c$ not necessarily nonnegative]
\label{rem:parabolic-max-principle-negative-c}
\uses{thm:parabolic-weak-maximum-principle-c-nonneg}
\dcref{ch7:7.1.4-rem-3}
\group{Maximum Principles}

Parabolic maximum principles remain valid for zeroth-order coefficients bounded below after an exponential change in time; nonnegativity of $c$ is not essential.
\end{remark}

\subsubsection{Harnack's inequality}

\begin{theorem}[Parabolic Harnack inequality]
\label{thm:parabolic-harnack-inequality}
\uses{thm:parabolic-weak-maximum-principle}
\dcref{ch7:7.1-thm-10}
\group{Maximum Principles}

Assume $u \in C^2_1(U_T)$ solves
\[
u_t + Lu = 0 \text{ in } U_T,
\tag{68}
\]
and
\[
u \geq 0 \text{ in } U_T.
\]
Suppose $V \subset\subset U$ is connected. Then for each $0 < t_1 < t_2 \leq T$, there exists a constant $C$ such that
\[
\sup_V u(\cdot, t_1) \leq C \inf_V u(\cdot, t_2).
\tag{69}
\]
The constant $C$ depends only on $V$, $t_1$, $t_2$, and the coefficients of $L$.
\end{theorem}

\begin{proof}
\sketch For smooth leading coefficients and no lower-order terms, a logarithmic/Bernstein computation with a parabolic cutoff yields a local gradient and time-oscillation estimate for positive solutions. Chaining overlapping interior cylinders gives the comparison between the earlier and later space-time regions.
\end{proof}

\subsubsection{Strong maximum principle}

\begin{theorem}[Strong maximum principle]
\label{thm:parabolic-strong-maximum-principle}
\uses{thm:parabolic-weak-maximum-principle, thm:parabolic-harnack-inequality}
\dcref{ch7:7.1-thm-11}
\group{Maximum Principles}

Assume $u \in C_1^2(U_T) \cap C(\overline{U_T})$ and
\[
c \equiv 0 \quad \text{in } U_T.
\]
Suppose also $U$ is connected.

(i) If
\[
u_t + Lu \leq 0 \quad \text{in } U_T
\]
and $u$ attains its maximum over $\overline{U_T}$ at a point $(x_0, t_0) \in U_T$, then
\[
u \text{ is constant on } U_{t_0}.
\]

(ii) Likewise, if
\[
u_t + Lu \geq 0 \quad \text{in } U_T
\]
and $u$ attains its minimum over $\overline{U_T}$ at a point $(x_0, t_0) \in U_T$, then
\[
u \text{ is constant on } U_{t_0}.
\]
\end{theorem}

\begin{proof}
\sketch Apply Harnack inequality along a chain of interior cylinders joining a point where the nonnegative solution vanishes to every point at an earlier time. The zero propagates through the chain, yielding constancy on the backward reachable region.
\end{proof}

\begin{remark}[Infinite propagation speed]
\label{rem:parabolic-infinite-propagation-speed}
\uses{thm:parabolic-strong-maximum-principle}
\dcref{ch7:7.1.4-rem-4}
\group{Maximum Principles}

A nontrivial nonnegative parabolic solution becomes positive throughout every connected spatial region at each later positive time; disturbances therefore have infinite propagation speed.
\end{remark}

\begin{theorem}[Strong maximum principle for $c \geq 0$]
\label{thm:parabolic-strong-maximum-principle-c-nonneg}
\uses{thm:parabolic-strong-maximum-principle, thm:parabolic-weak-maximum-principle-c-nonneg, thm:parabolic-harnack-inequality}
\dcref{ch7:7.1-thm-12}
\group{Maximum Principles}

Assume $u \in C_1^2(U_T) \cap C(\bar{U}_T)$ and
\[
c \geq 0 \quad \text{in } U_T.
\]
Suppose also $U$ is connected.

(i) If
\[
u_t + Lu \leq 0 \quad \text{in } U_T
\]
and $u$ attains a nonnegative maximum over $\bar{U}_T$ at a point $(x_0, t_0) \in U_T$, then
\[
u \text{ is constant on } U_{t_0}.
\]

(ii) Similarly, if
\[
u_t + Lu \geq 0 \quad \text{in } U_T
\]
and $u$ attains a nonpositive minimum over $\bar{U}_T$ at a point $(x_0, t_0) \in U_T$, then
\[
u \text{ is constant on } U_{t_0}.
\]
\end{theorem}

\begin{proof}
\sketch Apply the weak maximum principle to normalize the relevant sign and absorb the zeroth-order coefficient by an exponential time factor, then invoke the zero-order strong maximum principle.
\end{proof}

\section{Second-Order Hyperbolic Equations}
\label{sec:second-order-hyperbolic-equations}

\subsection{Definitions}

\subsubsection{Hyperbolic equations}

\begin{definition}[Uniform hyperbolicity]
\label{def:uniform-hyperbolicity}
\dcref{ch7:7.2.1-def-1}
\group{Hyperbolic PDE}
We say the partial differential operator $\frac{\partial^2}{\partial t^2} + L$ is (uniformly) hyperbolic if there exists a constant $\theta > 0$ such that
\[
\sum_{i,j=1}^{n} a^{ij}(x,t)\xi_i\xi_j \geq \theta|\xi|^2
\tag{4}
\]
for all $(x,t) \in U_T$, $\xi \in \mathbb{R}^n$.
\end{definition}

\subsubsection{Weak solutions}

\begin{remark}[Weak hyperbolic time derivative]
\label{rem:motivation-weak-solution-hyperbolic}
\uses{rem:motivation-weak-solution-parabolic}
\dcref{ch7:7.2.1-rem-1}
\group{Hyperbolic PDE}

Regard $u:[0,T]\to H_0^1(U)$ and $f:[0,T]\to L^2(U)$. Formal spatial integration by parts gives
\[
\langle u'',v\rangle+B[u,v;t]=(f,v) \qquad (v\in H_0^1(U)).
\tag{9}
\]
The equation implies
\[
u_{tt}=g^0+\sum_{j=1}^n g^j_{x_j},\qquad
g^0=f-\sum_{i=1}^n b^iu_{x_i}-cu,\quad
g^j=\sum_{i=1}^n a^{ij}u_{x_i},
\]
so $u''\in H^{-1}(U)$ is the natural weak time regularity.
\end{remark}

\begin{definition}[Weak solution of the hyperbolic initial/boundary-value problem]
\label{def:weak-solution-hyperbolic-bvp}
\uses{rem:motivation-weak-solution-hyperbolic}
\dcref{ch7:7.2.1-def-2}
\group{Hyperbolic PDE}

We say a function
\[
\mathbf{u} \in L^2(0, T; H^1_0(U)), \text{ with } \mathbf{u}' \in L^2(0, T; L^2(U)), \mathbf{u}'' \in L^2(0, T; H^{-1}(U)),
\]
is a weak solution of the \textit{hyperbolic initial/boundary-value problem} (1) provided

(i) $\langle \mathbf{u}'', v \rangle + B[\mathbf{u}, v; t] = (\mathbf{f}, v)$

for each $v \in H^1_0(U)$ and a.e. time $0 \le t \le T$, and

(ii) $\mathbf{u}(0) = g$, $\mathbf{u}'(0) = h$.
\end{definition}

\begin{remark}[Continuity of the weak solution]
\label{rem:weak-solution-hyperbolic-continuity}
\uses{def:weak-solution-hyperbolic-bvp, thm:calculus-abstract-space}
\dcref{ch7:7.2.1-rem-2}
\group{Hyperbolic PDE}

Every weak solution satisfies $u\in C([0,T];L^2(U))$ and $u'\in C([0,T];H^{-1}(U))$, so both initial conditions are meaningful.
\end{remark}

\subsection{Existence of Weak Solutions}

\subsubsection{Galerkin approximations}

\begin{theorem}[Construction of approximate solutions]
\label{thm:hyperbolic-galerkin-approximate-solutions}
\dcref{ch7:7.2-thm-1}
\group{Hyperbolic PDE}
For each integer $m = 1, 2, \ldots$, there exists a unique function $\mathbf{u}_m$ of the form (13) satisfying (14)--(16).
\end{theorem}

\begin{proof}
\sketch Expanding in the first $m$ basis functions converts the weak wave equation into a finite second-order linear ODE system with projected displacement and velocity data. Standard ODE theory gives the approximate solution.
\end{proof}

\subsubsection{Energy estimates}

\begin{theorem}[Energy estimates]
\label{thm:hyperbolic-energy-estimates}
\uses{thm:hyperbolic-galerkin-approximate-solutions}
\dcref{ch7:7.2-thm-2}
\group{Energy Estimates}

There exists a constant $C$, depending only on $U$, $T$ and the coefficients of $L$, such that
\[
\max_{0 \le t \le T} \left( \|u_m(t)\|_{H^1_0(U)} + \|u'_m(t)\|_{L^2(U)} \right) + \|u''_m\|_{L^2(0,T;H^{-1}(U))}
\]
\[
\le C \left( \|\mathbf{f}\|_{L^2(0,T;L^2(U))} + \|g\|_{H^1_0(U)} + \|h\|_{L^2(U)} \right),
\tag{19}
\]
for $m = 1, 2, \ldots$
\end{theorem}

\begin{proof}
\sketch Multiply the Galerkin equations by the coefficient velocities and sum. Ellipticity controls the spatial energy, time differentiation of the coefficients contributes bounded error terms, and Cauchy inequality plus Gronwall yields the displacement, velocity, and acceleration bounds.
\end{proof}

\subsubsection{Existence and uniqueness}

\begin{theorem}[Existence of weak solution]
\label{thm:hyperbolic-existence-weak-solution}
\uses{thm:hyperbolic-galerkin-approximate-solutions, thm:hyperbolic-energy-estimates, def:weak-solution-hyperbolic-bvp}
\dcref{ch7:7.2-thm-3}
\group{Hyperbolic PDE}

There exists a weak solution of (1).
\end{theorem}

\begin{proof}
\sketch Use the energy bounds to extract weakly convergent subsequences for $u_m$, $u_m^{\prime}$, and $u_m^{\prime\prime}$. Pass to the Galerkin identity by density, identify both initial conditions by time integration by parts, and use abstract continuity in time.
\end{proof}

\begin{remark}[Improved regularity of the weak solution]
\label{rem:hyperbolic-existence-remark}
\uses{thm:hyperbolic-existence-weak-solution, thm:hyperbolic-energy-estimates}
\dcref{ch7:7.2.2-rem-1}
\group{Hyperbolic PDE}

The energy estimates sharpen the weak-solution regularity to
\[
u\in L^\infty(0,T;H_0^1(U)),\qquad
u'\in L^\infty(0,T;L^2(U)),\qquad
u''\in L^2(0,T;H^{-1}(U)).
\]
\end{remark}

\begin{theorem}[Uniqueness of weak solution]
\label{thm:hyperbolic-uniqueness-weak-solution}
\uses{def:weak-solution-hyperbolic-bvp}
\dcref{ch7:7.2-thm-4}
\group{Hyperbolic PDE}

A weak solution of (1) is unique.
\end{theorem}

\begin{proof}
\sketch For the difference $u$, use the backward primitive $v(t)=\int_t^s u(\tau)\,d\tau$ as a time-dependent test function. Integration in time produces an energy identity at $s$; ellipticity and Gronwall imply $u=0$.
\end{proof}

\begin{remark}[Admissible test functions for uniqueness]
\label{rem:hyperbolic-uniqueness-remark}
\uses{thm:hyperbolic-uniqueness-weak-solution}
\dcref{ch7:7.2.2-rem-2}
\group{Hyperbolic PDE}

The available regularity does not permit using $u'(t)$ directly as a test function; the uniqueness proof instead uses a time primitive of $u$.
\end{remark}

\subsection{Regularity}

\begin{remark}[Formal hyperbolic regularity estimates]
\label{rem:hyperbolic-regularity-motivation}
\uses{thm:wave-equation-nonhomogeneous-duhamel}
\dcref{ch7:7.2.3-rem-1}
\group{Regularity}

For a smooth rapidly decaying solution of $u_{tt}-\Delta u=f$ on $\mathbb R^n\times(0,T)$ with data $(g,h)$, the wave-energy identity and Gronwall inequality give
\[
\sup_{0\le t\le T}\int_{\mathbb R^n}|Du|^2+u_t^2\,dx
\le C\left(\int_0^T\int_{\mathbb R^n}f^2\,dx\,dt+\int_{\mathbb R^n}|Dg|^2+h^2\,dx\right).
\tag{42}
\]
Applying this estimate to $u_t$ and using $-\Delta u=f-u_{tt}$ yields
\[
\sup_{0\le t\le T}\int_{\mathbb R^n}|D^2u|^2+|Du_t|^2+u_{tt}^2\,dx
\le C\left(\int_0^T\int_{\mathbb R^n}f_t^2+f^2\,dx\,dt+\int_{\mathbb R^n}|D^2g|^2+|Dh|^2\,dx\right).
\tag{47}
\]
\end{remark}

\begin{theorem}[Improved regularity]
\label{thm:hyperbolic-improved-regularity}
\uses{def:weak-solution-hyperbolic-bvp, thm:parabolic-improved-regularity}
\dcref{ch7:7.2-thm-5}
\group{Regularity}

(i) Assume
\[
g \in H_0^1(U), \quad h \in L^2(U), \quad f \in L^2(0,T;L^2(U)),
\]
and suppose also $u \in L^2(0,T;H_0^1(U))$, with $u' \in L^2(0,T;L^2(U))$, $u'' \in L^2(0,T;H^{-1}(U))$, is the weak solution of the problem
\[
\begin{cases} u_{tt} + Lu = f & \text{in } U_T \\ u = 0 & \text{on } \partial U \times [0,T] \\ u = g, u_t = h & \text{on } U \times \{t = 0\}. \end{cases}
\tag{49}
\]
Then in fact
\[
u \in L^\infty(0,T;H_0^1(U)), \quad u' \in L^\infty(0,T;L^2(U)),
\]
and we have the estimate
\[
\operatorname{ess\,sup}_{0 \le t \le T} \left( \|u(t)\|_{H_0^1(U)} + \|u'(t)\|_{L^2(U)} \right) \le C\left( \|f\|_{L^2(0,T;L^2(U))} + \|g\|_{H_0^1(U)} + \|h\|_{L^2(U)} \right).
\tag{50}
\]

(ii) If, in addition,
\[
g \in H^2(U), \quad h \in H_0^1(U), \quad f' \in L^2(0,T;L^2(U)),
\]
then
\[
u \in L^\infty(0,T;H^2(U)), \quad u' \in L^\infty(0,T;H_0^1(U)),
\]
\[
u'' \in L^\infty(0,T;L^2(U)), \quad u''' \in L^2(0,T;H^{-1}(U)),
\]
with the estimate:
\[
\operatorname{ess\,sup}_{0 \le t \le T} \left( \|u(t)\|_{H^2(U)} + \|u'(t)\|_{H_0^1(U)} + \|u''(t)\|_{L^2(U)} \right) + \|u'''\|_{L^2(0,T;H^{-1}(U))} \le C\left( \|f\|_{H^1(0,T;L^2(U))} + \|g\|_{H^2(U)} + \|h\|_{H^1(U)} \right).
\tag{51}
\]
\end{theorem}

\begin{proof}
\sketch Differentiate the Galerkin system and estimate the differentiated hyperbolic energy, using the compatibility equation for the initial acceleration. Passage to the limit gives higher time regularity, and elliptic boundary regularity applied to $Lu=f-u^{\prime\prime}$ gives $H^2$ spatial regularity.
\end{proof}

\begin{remark}[Precise versions of the formal estimates]
\label{rem:hyperbolic-improved-regularity-remark}
\uses{thm:hyperbolic-improved-regularity, rem:hyperbolic-regularity-motivation}
\dcref{ch7:7.2.3-rem-2}
\group{Regularity}

Parts (i) and (ii) of \cref{thm:hyperbolic-improved-regularity} are the bounded-domain, variable-coefficient versions of (42) and (47).
\end{remark}

\begin{remark}[Symmetric operators and eigenfunction bases]
\label{rem:hyperbolic-symmetric-eigenfunctions}
\uses{thm:hyperbolic-improved-regularity}
\dcref{ch7:7.2.3-rem-3}
\group{Regularity}

If $L$ is symmetric, the Galerkin basis may be chosen as an $H_0^1(U)$ eigenfunction basis of $L$, eliminating estimate (59).
\end{remark}

\begin{theorem}[Higher regularity]
\label{thm:hyperbolic-higher-regularity}
\uses{thm:hyperbolic-improved-regularity, thm:mappings-into-better-spaces}
\dcref{ch7:7.2-thm-6}
\group{Regularity}

Assume
\[
\begin{cases}
g \in H^{m+1}(U), \quad h \in H^m(U), \\
\dfrac{d^k f}{dt^k} \in L^2(0, T; H^{m-k}(U)) \quad (k = 0, \ldots, m).
\end{cases}
\]
Suppose also the following $m^{\text{th}}$-order compatibility conditions hold:
\[
\begin{cases}
g_0 := g \in H^1_0(U), \quad h_1 := h \in H^1_0(U), \ldots, \\
g_{2l} := \dfrac{d^{2l}f}{dt^{2l}}(\cdot, 0) - Lg_{2l-2} \in H^1_0(U)
  \quad \text{(if } m = 2l\text{)}, \\
h_{2l+1} := \dfrac{d^{2l+1}f}{dt^{2l+1}}(\cdot, 0) - Lh_{2l-1} \in H^1_0(U)
  \quad \text{(if } m = 2l + 1\text{)}.
\end{cases}
\tag{62}
\]
Then
\[
\frac{d^k u}{dt^k} \in L^\infty(0, T; H^{m+1-k}(U)) \quad (k = 0, \ldots, m+1),
\tag{63}
\]
and we have the estimate
\[
\operatorname{ess\,sup}_{0 \le t \le T} \sum_{k=0}^{m+1} \left\| \frac{d^k u}{dt^k} \right\|_{H^{m+1-k}(U)}
\]
\[
\le C \left( \sum_{k=0}^{m} \left\| \frac{d^k f}{dt^k} \right\|_{L^2(0,T;H^{m-k}(U))} + \|g\|_{H^{m+1}(U)} + \|h\|_{H^m(U)} \right).
\tag{64}
\]
\end{theorem}

\begin{proof}
\sketch Induct on time derivatives, applying the improved-regularity estimate to each differentiated equation and using the recursive compatibility conditions. Higher elliptic regularity recovers the corresponding spatial derivatives.
\end{proof}

\begin{remark}[Compatibility conditions]
\label{rem:hyperbolic-higher-regularity-remark}
\uses{thm:hyperbolic-higher-regularity, thm:mappings-into-better-spaces}
\dcref{ch7:7.2.3-rem-4}
\group{Regularity}

The time traces obey
\[
f(0)\in H^{m-1}(U),\quad f'(0)\in H^{m-2}(U),\quad\ldots,\quad f^{(m-2)}(0)\in H^1(U),
\tag{65}
\]
and the recursively defined data satisfy
\[
g_0\in H^{m+1}(U),\ h_1\in H^m(U),\ g_2\in H^{m-1}(U),\ h_3\in H^{m-2}(U),\ldots.
\tag{66}
\]
Compatibility requires each listed datum to have zero trace on $\partial U$.
\end{remark}

\begin{theorem}[Infinite differentiability]
\label{thm:hyperbolic-infinite-differentiability}
\uses{thm:hyperbolic-higher-regularity}
\dcref{ch7:7.2-thm-7}
\group{Regularity}

Assume
\[
g, h \in C^\infty(\bar{U}), \quad f \in C^\infty(\bar{U}_T),
\]
and the $m^{\text{th}}$-order compatibility conditions hold for $m = 0, 1, \ldots$.

Then the hyperbolic initial/boundary-value problem (1) has a unique solution
\[
u \in C^\infty(\bar{U}_T).
\]
\end{theorem}

\begin{proof}
\sketch Higher hyperbolic regularity gives Sobolev estimates of arbitrary order; Sobolev embedding in space and time yields smoothness.
\end{proof}

\subsection{Propagation of Disturbances}

\begin{theorem}[Finite propagation speed]
\label{thm:hyperbolic-finite-propagation-speed}
\uses{def:uniform-hyperbolicity, thm:wave-equation-finite-propagation-speed}
\dcref{ch7:7.2-thm-8}
\group{Hyperbolic PDE}

Assume $u$ is a smooth solution of the hyperbolic equation (72). If $u \equiv u_t \equiv 0$ on $C_0$, then $u \equiv 0$ within the cone $C$.
\end{theorem}

\begin{proof}
\sketch Integrate the local energy over balls shrinking at the characteristic speed. Differentiating this cone energy gives an interior contribution controlled by the energy and a nonpositive lateral-boundary flux by hyperbolicity; Gronwall and zero initial energy imply vanishing in the cone.
\end{proof}

\begin{remark}[Finite domain of dependence]
\label{rem:hyperbolic-finite-propagation-remark}
\uses{thm:hyperbolic-finite-propagation-speed}
\dcref{ch7:7.2.4-rem-1}
\group{Hyperbolic PDE}

For initial data $u(0)=g$ and $u_t(0)=h$, the value $u(x_0,t_0)$ depends only on the restrictions of $g$ and $h$ to the base $C_0$ of the backward characteristic cone.
\end{remark}

\subsection{Equations in Two Variables}

\section{Hyperbolic Systems of First-Order Equations}
\label{sec:hyperbolic-systems-first-order-equations}

\subsection{Definitions}

\begin{definition}[Hyperbolic system]
\label{def:hyperbolic-system}
\dcref{ch7:7.3.1-def-1}
\group{Hyperbolic Systems}
The system of PDE (1) is called hyperbolic if the $m \times m$ matrix $\mathbf{B}(x, t; y)$ is diagonalizable for each $x, y \in \mathbb{R}^n$, $t \geq 0$.
\end{definition}

\begin{definition}[Symmetric and strictly hyperbolic systems]
\label{def:symmetric-strictly-hyperbolic-system}
\uses{def:hyperbolic-system}
\dcref{ch7:7.3.1-def-2}
\group{Hyperbolic Systems}

(i) We say (1) is a symmetric hyperbolic system if $B_j(x, t)$ is a symmetric $m \times m$ matrix for each $x \in \mathbb{R}^n$, $t \geq 0$ ($j = 1, \ldots, m$).

(ii) The system (1) is strictly hyperbolic if for each $x, y \in \mathbb{R}^n$, $y \neq 0$, and each $t \geq 0$, the matrix $\mathbf{B}(x, t; y)$ has $m$ distinct real eigenvalues:
\[
\lambda_1(x, t; y) < \lambda_2(x, t; y) < \cdots < \lambda_m(x, t; y).
\]
\end{definition}

\begin{remark}[Plane-wave characterization of hyperbolicity]
\label{rem:motivation-hyperbolicity-plane-wave}
\uses{def:hyperbolic-system, def:traveling-wave-plane-wave}
\dcref{ch7:7.3.1-rem-1}
\group{Hyperbolic Systems}

For constant matrices and $\mathbf f=0$, a plane wave
\[
\mathbf u(x,t)=\mathbf v(y\cdot x-\sigma t)
\tag{4}
\]
satisfies
\[
\left(-\sigma I+\sum_{j=1}^ny_jB_j\right)\mathbf v'=0.
\]
Thus $\sigma$ is an eigenvalue of $\mathbf B(y)=\sum_jy_jB_j$ and $\mathbf v'$ an eigenvector. Hyperbolicity supplies $m$ independent real plane-wave modes; for $|y|=1$, the eigenvalues are their wave speeds.
\end{remark}

\subsection{Symmetric Hyperbolic Systems}

\begin{remark}[Symmetric systems generalize second-order hyperbolic PDE]
\label{rem:symmetric-hyperbolic-general-systems}
\uses{def:symmetric-strictly-hyperbolic-system, def:uniform-hyperbolicity}
\dcref{ch7:7.3.2-rem-1}
\group{Hyperbolic Systems}

A system
\[
B_0\mathbf u_t+\sum_{j=1}^nB_j\mathbf u_{x_j}=\mathbf f
\tag{9}
\]
is symmetric when every $B_j$ is symmetric and $B_0$ is positive definite. The scalar equation
\[
v_{tt}-\sum_{i,j=1}^na^{ij}v_{x_ix_j}=0
\tag{10}
\]
has this form for $\mathbf u=(v_{x_1},\ldots,v_{x_n},v_t)$, with $B_0$ positive definite by uniform hyperbolicity.
\end{remark}

\subsubsection{Weak solutions}

\begin{definition}[Weak solution of a symmetric hyperbolic system]
\label{def:weak-solution-symmetric-hyperbolic-system}
\uses{rem:symmetric-hyperbolic-general-systems}
\dcref{ch7:7.3.2-def-1}
\group{Hyperbolic Systems}

We say
\[
\mathbf{u} \in L^2(0,T; H^1(\mathbb{R}^n; \mathbb{R}^m)), \text{ with } \mathbf{u}' \in L^2(0,T; L^2(\mathbb{R}^n; \mathbb{R}^m)),
\]
is a weak solution of the initial-value problem (5) for the symmetric hyperbolic system provided

(i) $(\mathbf{u}',\mathbf{v}) + B[\mathbf{u},\mathbf{v};t] = (\mathbf{f},\mathbf{v})$

for each $\mathbf{v} \in H^1(\mathbb{R}^n; \mathbb{R}^m)$ and a.e. $0 \le t \le T$, and

(ii) $\mathbf{u}(0) = \mathbf{g}$.

Here and afterwards $(\cdot, \cdot)$ denotes the inner product in $L^2(\mathbb{R}^n; \mathbb{R}^m)$.
\end{definition}

\begin{remark}[Continuity of the weak solution]
\label{rem:weak-solution-symmetric-hyperbolic-continuity}
\uses{def:weak-solution-symmetric-hyperbolic-system, thm:calculus-abstract-space}
\dcref{ch7:7.3.2-rem-2}
\group{Hyperbolic Systems}

Every weak solution belongs to $C([0,T];L^2(\mathbb R^n;\mathbb R^m))$, so the initial condition is meaningful.
\end{remark}

\subsubsection{Vanishing viscosity method}

\begin{theorem}[Existence of approximate solutions]
\label{thm:symmetric-hyperbolic-approximate-solutions}
\uses{def:fundamental-solution-heat-equation}
\dcref{ch7:7.3-thm-1}
\group{Hyperbolic Systems}

For each $\epsilon > 0$, there exists a unique solution $\mathbf{u}^\epsilon$ of (11), with
\[
\mathbf{u}^\epsilon \in L^2(0,T; H^3(\mathbb{R}^n; \mathbb{R}^m)), \quad (\mathbf{u}^\epsilon)' \in L^2(0,T; H^1(\mathbb{R}^n; \mathbb{R}^m)).
\tag{12}
\]
\end{theorem}

\begin{proof}
\sketch Add the viscosity term $-\varepsilon\Delta u^{\varepsilon}$ with compatible boundary data. The resulting uniformly parabolic system has a smooth solution by the parabolic existence theory.
\end{proof}

\subsubsection{Energy estimates}

\begin{theorem}[Energy estimates]
\label{thm:symmetric-hyperbolic-energy-estimates}
\uses{thm:symmetric-hyperbolic-approximate-solutions}
\dcref{ch7:7.3-thm-2}
\group{Energy Estimates}

There exists a constant $C$, depending only on $n$ and the coefficients, such that
\[
\max_{0 \leq t \leq T} \|u^\epsilon(t)\|_{H^1(\mathbb{R}^n;\mathbb{R}^m)} + \|u_t^\epsilon\|_{L^2(0,T;L^2(\mathbb{R}^n;\mathbb{R}^m))}
\]
\[
\leq C(\|g\|_{H^1(\mathbb{R}^n;\mathbb{R}^m)} + \|f\|_{L^2(0,T;H^1(\mathbb{R}^n;\mathbb{R}^m))} + \|f'\|_{L^2(0,T;L^2(\mathbb{R}^n;\mathbb{R}^m))})
\tag{17}
\]
for each $0 < \epsilon \leq 1$.
\end{theorem}

\begin{proof}
\sketch Take the Euclidean inner product with $u^{\varepsilon}$ and integrate over space. Symmetry permits integration by parts of the principal matrices, the viscosity term is dissipative, and bounded lower-order terms are controlled by Gronwall uniformly in $\varepsilon$.
\end{proof}

\subsubsection{Existence and uniqueness}

\begin{theorem}[Existence of weak solution]
\label{thm:symmetric-hyperbolic-existence-weak-solution}
\uses{thm:symmetric-hyperbolic-energy-estimates, def:weak-solution-symmetric-hyperbolic-system}
\dcref{ch7:7.3-thm-3}
\group{Hyperbolic Systems}

There exists a weak solution of the initial value problem (5).
\end{theorem}

\begin{proof}
\sketch The uniform viscosity estimates yield a weakly convergent subsequence. The viscous term tends to zero, so passage to the limit in the weak formulation gives a solution; time integration by parts recovers the initial data and continuity.
\end{proof}

\begin{theorem}[Uniqueness of weak solution]
\label{thm:symmetric-hyperbolic-uniqueness-weak-solution}
\uses{def:weak-solution-symmetric-hyperbolic-system}
\dcref{ch7:7.3-thm-4}
\group{Hyperbolic Systems}

A weak solution of (5) is unique.
\end{theorem}

\begin{proof}
\sketch Apply the symmetric-system energy identity to the difference of two solutions and use Gronwall.
\end{proof}

\subsection{Systems with Constant Coefficients}

\begin{theorem}[Existence of solution]
\label{thm:constant-coefficient-hyperbolic-existence}
\uses{def:fractional-sobolev-space, def:fourier-transform-l1}
\dcref{ch7:7.3-thm-5}
\group{Hyperbolic Systems}

Assume
\[
\mathbf{g} \in H^s(\mathbb{R}^n; \mathbb{R}^m) \quad \left(s > \frac{n}{2} + m\right).
\]
Then there is a unique solution $\mathbf{u} \in C^1([0, \infty); \mathbb{R}^m)$ of the initial-value problem (32), (33).
\end{theorem}

\begin{proof}
\sketch Fourier transformation changes the PDE into an ODE $\hat u_t+iA(\xi)\hat u=\hat f$. Hyperbolicity diagonalizes $A(\xi)$ with real spectrum and uniform estimates; solve the ODE, bound the weighted Fourier norm, and invert the transform.
\end{proof}

\section{Semigroup Theory}
\label{sec:semigroup-theory}

\subsection{Definitions, Elementary Properties}

\subsubsection{Semigroups}

\begin{definition}[Semigroup, contraction semigroup]
\label{def:semigroup-contraction-semigroup}
\dcref{ch7:7.4.1-def-1}
\group{Semigroup Theory}
(i) A family $\{S(t)\}_{t \geq 0}$ of bounded linear operators mapping $X$ into $X$ is called a semigroup if conditions (4)--(6) are satisfied.

(ii) We say $\{S(t)\}_{t \geq 0}$ is a contraction semigroup if in addition
\[
\|S(t)\| \leq 1 \quad (t \geq 0),
\tag{7}
\]
$\|\cdot\|$ here denoting the operator norm. Thus
\[
\|S(t)u\| \leq \|u\| \quad (t \geq 0, u \in X).
\]
\end{definition}

\subsubsection{Elementary properties, generators}

\begin{definition}[Infinitesimal generator]
\label{def:semigroup-generator}
\uses{def:semigroup-contraction-semigroup}
\dcref{ch7:7.4.1-def-2}
\group{Semigroup Theory}

Write
\[
D(A) := \left\{u \in X \mid \lim_{t \to 0+} \frac{S(t)u - u}{t} \text{ exists in } X\right\}
\tag{8}
\]
and
\[
Au := \lim_{t \to 0+} \frac{S(t)u - u}{t} \quad (u \in D(A)).
\tag{9}
\]
We call $A : D(A) \to X$ the (infinitesimal) generator of the semigroup $\{S(t)\}_{t \geq 0}$; $D(A)$ is the domain of $A$.
\end{definition}

\begin{theorem}[Differential properties of semigroups]
\label{thm:semigroup-differential-properties}
\uses{def:semigroup-generator}
\dcref{ch7:7.4-thm-1}
\group{Semigroup Theory}

Assume $u \in D(A)$. Then
(i) $S(t)u \in D(A)$ for each $t \geq 0$,
(ii) $AS(t)u = S(t)Au$ for each $t > 0$,
(iii) the mapping $t \mapsto S(t)u$ is differentiable for each $t > 0$,
and
(iv) $\frac{d}{dt}S(t)u = AS(t)u$ \quad $(t > 0)$.
\end{theorem}

\begin{proof}
\sketch Use the semigroup law to rewrite time difference quotients of $S(t)u$ as $S(t)$ applied to the generator difference quotient. Strong continuity gives differentiability and invariance of the generator domain, and integration yields the fundamental identity.
\end{proof}

\begin{remark}[Continuous differentiability of the flow]
\label{rem:semigroup-c1-remark}
\uses{thm:semigroup-differential-properties}
\dcref{ch7:7.4.1-rem-1}
\group{Semigroup Theory}

If $u\in D(A)$, then $t\mapsto S(t)u$ is $C^1$ on $(0,\infty)$ because $AS(t)u=S(t)Au$ is continuous.
\end{remark}

\begin{theorem}[Properties of generators]
\label{thm:semigroup-generator-properties}
\uses{def:semigroup-generator, thm:semigroup-differential-properties}
\dcref{ch7:7.4-thm-2}
\group{Semigroup Theory}

(i) The domain $D(A)$ is dense in $X$,

and

(ii) $A$ is a closed operator.
\end{theorem}

\begin{proof}
\sketch Time averages $t^{-1}\int_0^tS(s)u\,ds$ lie in the generator domain and converge to $u$, proving density. Passing limits through the defining difference quotients proves closedness, while the semigroup differential identity gives invariance and commutation.
\end{proof}

\begin{remark}[Meaning of ``closed operator'']
\label{rem:closed-operator-definition}
\uses{thm:semigroup-generator-properties}
\dcref{ch7:7.4.1-rem-2}
\group{Semigroup Theory}

Closedness means that $u_k\in D(A)$, $u_k\to u$, and $Au_k\to v$ imply
\[
u\in D(A),\qquad Au=v.
\]
\end{remark}

\subsubsection{Resolvents}

\begin{definition}[Resolvent set and resolvent operator]
\label{def:resolvent-set-operator}
\uses{def:semigroup-generator}
\dcref{ch7:7.4.1-def-3}
\group{Semigroup Theory}

(i) We say a real number $\lambda$ belongs to $\rho(A)$, the resolvent set of $A$, provided the operator
\[
\lambda I - A : D(A) \to X
\]
is one-to-one and onto.

(ii) If $\lambda \in \rho(A)$, the resolvent operator $R_\lambda : X \to X$ is defined by
\[
R_\lambda u := (\lambda I - A)^{-1}u.
\]
\end{definition}

\begin{theorem}[Properties of resolvent operators]
\label{thm:resolvent-operator-properties}
\uses{def:resolvent-set-operator}
\dcref{ch7:7.4-thm-3}
\group{Semigroup Theory}

(i) If $\lambda, \mu \in \rho(A)$, we have
\[
R_\lambda - R_\mu = (\mu - \lambda)R_\lambda R_\mu \quad \text{(resolvent identity)}
\tag{12}
\]
and
\[
R_\lambda R_\mu = R_\mu R_\lambda.
\tag{13}
\]

(ii) If $\lambda > 0$, then $\lambda \in \rho(A)$,
\[
R_\lambda u = \int_0^\infty e^{-\lambda t}S(t)u \, dt \quad (u \in X),
\tag{14}
\]
and so $\|R_\lambda\| \leq \frac{1}{\lambda}$.
\end{theorem}

\begin{proof}
\sketch For $\lambda>0$, define $R_\lambda u=\int_0^\infty e^{-\lambda t}S(t)u\,dt$. Integration by parts and the generator identity show $(\lambda I-A)R_\lambda=I=R_\lambda(\lambda I-A)$; contraction gives $\|R_\lambda\|\le\lambda^{-1}$.
\end{proof}

\begin{remark}[Resolvent as Laplace transform]
\label{rem:resolvent-laplace-transform-semigroup}
\uses{thm:resolvent-operator-properties, ex:resolvent-laplace-transform}
\dcref{ch7:7.4.1-rem-3}
\group{Semigroup Theory}

For $\lambda>0$, the resolvent is the Laplace transform
\[
R_\lambda u=\int_0^\infty e^{-\lambda t}S(t)u\,dt.
\]
\end{remark}

\subsection{Generating Contraction Semigroups}
\label{sec:generating-contraction-semigroups}

\begin{theorem}[Hille--Yosida Theorem]
\label{thm:hille-yosida-theorem}
\uses{def:semigroup-contraction-semigroup, def:semigroup-generator, thm:resolvent-operator-properties}
\dcref{ch7:7.4-thm-4}
\group{Semigroup Theory}

Let $A$ be a closed, densely-defined linear operator on $X$. Then $A$ is the generator of a contraction semigroup $\{S(t)\}_{t \geq 0}$ if and only if
\[
(0,\infty) \subset \rho(A) \text{ and } \|R_\lambda\| \leq \frac{1}{\lambda} \text{ for } \lambda > 0.
\tag{17}
\]
\end{theorem}

\begin{proof}
\sketch Necessity follows from the Laplace-transform resolvent. Conversely define the bounded Yosida approximants $A_\lambda=\lambda AR_\lambda$ and semigroups $e^{tA_\lambda}$. Resolvent estimates give contraction and Cauchy bounds; their strong limit is a contraction semigroup whose generator is the closure of $A$.
\end{proof}

\begin{remark}[$\omega$-contractive semigroups]
\label{rem:omega-contractive-semigroup}
\uses{thm:hille-yosida-theorem}
\dcref{ch7:7.4.2-rem-1}
\group{Semigroup Theory}

A semigroup is $\omega$-contractive when $\|S(t)\|\le e^{\omega t}$. A closed densely defined operator $A$ generates such a semigroup iff
\[
(\omega,\infty)\subset\rho(A),\qquad
\|R_\lambda\|\le\frac1{\lambda-\omega}\quad(\lambda>\omega).
\tag{23}
\]
\end{remark}

\subsection{Applications}

\subsubsection{Second-order parabolic PDE}

\begin{theorem}[Second-order parabolic PDE as semigroups]
\label{thm:parabolic-pde-as-semigroup}
\uses{thm:hille-yosida-theorem, thm:energy-estimates, thm:boundary-h2-regularity, thm:first-existence-weak-solutions}
\dcref{ch7:7.4-thm-5}
\group{Semigroup Theory}

The operator $A$ generates a $\gamma$-contraction semigroup $\{S(t)\}_{t \ge 0}$ on $L^2(U)$.
\end{theorem}

\begin{proof}
\sketch On $L^2(U)$ take $Au=-Lu$ with domain $H^2(U)\cap H_0^1(U)$. Elliptic energy estimates give dissipativity and elliptic solvability gives the resolvent range condition, so Hille--Yosida generates a contraction semigroup; its orbit is the parabolic solution.
\end{proof}

\begin{remark}[Comparison with the Galerkin method]
\label{rem:parabolic-semigroup-remark}
\uses{thm:parabolic-pde-as-semigroup, thm:parabolic-existence-weak-solution}
\dcref{ch7:7.4.3-rem-1}
\group{Semigroup Theory}

The semigroup construction requires time-independent coefficients and directly yields a more regular solution. The Galerkin construction permits time-dependent coefficients.
\end{remark}

\subsubsection{Second-order hyperbolic PDE}

\begin{theorem}[Second-order hyperbolic PDE as semigroups]
\label{thm:hyperbolic-pde-as-semigroup}
\uses{thm:hille-yosida-theorem, thm:parabolic-pde-as-semigroup}
\dcref{ch7:7.4-thm-6}
\group{Semigroup Theory}

The operator $A$ generates a contraction semigroup $\{S(t)\}_{t \geq 0}$ on $H_0^1(U) \times L^2(U)$.
\end{theorem}

\begin{proof}
\sketch Rewrite the wave equation as a first-order system on $H_0^1(U)\times L^2(U)$ with the energy inner product. The block operator is dissipative and its resolvent reduces to the elliptic resolvent, so Hille--Yosida yields the hyperbolic evolution.
\end{proof}

\section{References}
\label{sec:chapter7-references}
\textbf{Parabolic regularity.} \cite{Krylov1996}, \cite{LadyzhenskayaSolonnikovUraltseva1968}, \cite{Lieberman1996}.
\textbf{Hyperbolic equations.} \cite{Ladyzhenskaya1985}, \cite{Lions1961}, \cite{LionsMagenes1972}, \cite{Wloka1987}.
\textbf{Hyperbolic systems.} \cite{John1982}, \cite{Treves1975}.
\textbf{Semigroups.} \cite{Friedman1969}, \cite{Yosida1980}.
\textbf{Further results.} \cite{BrezisEvans1979}, \cite{LadyzhenskayaUraltseva1968}, \cite{Lieberman1996}, \cite{Friedman1969}, \cite{Yosida1980}.
